% Peak extraction protocol -- append to the comparative study setup.
% Requires: amsmath, booktabs.

\subsection{Peak extraction and characterisation}
\label{subsec:peak-extraction}

Each detector $m \in \mathcal{M}$ produces a score series $s^{(m)}_t$ over a
common time index. No anomaly labels are used. Because detector score scales
differ, we define a comparable operating point using the same empirical
quantile $\alpha$ for every detector:
\begin{equation}
  u_m = Q_{\alpha}\!\left(s^{(m)}_1,\ldots,s^{(m)}_T\right),
  \qquad I^{(m)}_t = \mathbf{1}\{s^{(m)}_t > u_m\}.
  \label{eq:global-percentile}
\end{equation}
Thus each detector has approximately the same raw exceedance rate
$1-\alpha$. Ties can reduce the realised rate and are therefore reported.

\paragraph{Runs declustering.}
Raw exceedances separated by at most $r$ observations are assigned to one
cluster, represented by its largest score. The resulting event set is
$\mathcal{E}^{(m)}$. We report both the raw exceedance count $N_m$ and event
count $K_m$. Their ratio
\begin{equation}
  \hat{\theta}_m = K_m/N_m
  \label{eq:extremal-index}
\end{equation}
characterises persistence: smaller values indicate that extreme scores arrive
in longer or more closely spaced bursts.

\paragraph{Seasonal diagnostic.}
Seasonality is measured rather than removed. Let $c_m(d)$ be the mean score for
calendar day $d$ across years. We report the fraction of score variance
explained by this calendar-day climatology,
\begin{equation}
  S_m = 1 -
  \frac{\sum_t (s^{(m)}_t-c_m(d_t))^2}
       {\sum_t (s^{(m)}_t-\bar{s}^{(m)})^2},
  \label{eq:seasonal-strength}
\end{equation}
clipped to $[0,1]$. Month--day keys keep the climatology aligned through leap
years. A value near zero indicates little repeatable annual structure in the
scores; a value near one indicates a strongly repeating calendar-day pattern.

A month-of-year event histogram is retained as a direct visual diagnostic of
seasonal concentration in the extracted events.

\begin{table}[t]
  \centering
  \caption{Peak extraction and characterisation parameters.}
  \label{tab:peak-params}
  \begin{tabular}{llc}
    \toprule
    Symbol & Meaning & Value \\
    \midrule
    $\alpha$ & global extraction percentile & $0.99$ \\
    $r$      & declustering and event-matching window & $15$ days \\
    $\kappa$ & minimum detectors for consensus & $3$ \\
    \bottomrule
  \end{tabular}
\end{table}

\paragraph{Cross-detector agreement.}
For each pair of detectors, raw exceedance indicators before declustering are
compared without temporal widening, over the span covered by both score series:
\begin{equation}
  J(m,m') =
  \frac{|\mathcal{I}^{(m)} \cap \mathcal{I}^{(m')}|}
       {|\mathcal{I}^{(m)} \cup \mathcal{I}^{(m')}|}.
\end{equation}
Consensus events are maximal intervals supported by at least $\kappa$
detectors. Concordance between detectors with different inductive biases is a
useful prioritisation signal, but does not establish physical truth.

\paragraph{Outputs and limitations.}
For every series the analysis writes the per-detector diagnostic table, monthly
histogram, detector Jaccard matrix, consensus intervals, event masks, and the
existing detailed interpretation figures. Those figures centre on each
declustered cluster maximum and separately mark the other raw abnormal points
in its runs cluster. For subsequence detectors, the cluster maximum selects the
complete subsequence shown with its nearest training neighbour; MP comparisons
are z-normalized. An ``event'' means only that a score
is extreme relative to that detector's own distribution. Percentile extraction
does not estimate a false-alarm probability, correct score drift, or prove that
a physical anomaly occurred. Sensitivity analyses vary $\alpha$ and $r$.
